From Semtech's course:

- There are many wireless technologies: Bluetooth, Cellular, Wi-Fi, LoRaWAN, Zigbee, NB-IoT, and others
- 3 factors to keep in mind: range, bandwidth/datarate and power consumption
- LPWAN:
    - Low power (years on a battery or autonomous using solar panels)
    - Long range: 2-5km in urban and >10km in rural areas
- There are 3 LPWAN technologies: Sigfox, Narrowband-IoT (NB-IoT) and LoRa 
- LoRa datarate: 300bps-11kbps
- Max packet size in Europe: 222 bytes @ high data rate and 51 bytes @ low data rate
- Devices can send anytime, since gateways are alway listening
- Gateways can send during devices' RX windows
- End devices can be put into sleep mode to conserve energy
- LoRa can also be used for geolocation, albeit not as precise as GPS modules
- LoRa vs LoRaWAN:
    - LoRa is modulation technology
    - LoRaWAN is the standard that defines everything related to packets and maintained by LoRa Alliance
- LoRa uses ISM band (industrial, scientific and medical), 433 and 868MHz in Europe
- LoRaWan network comprises: application, join and network server
- 10 channels in total, of which 3 are fixed, i.e have to be supported by all devices
- EU 868 duty cycle limitations: 0.1, 1 or 10\% (for gateways only)
- Regulations differ depending on the region -> devices that work in the EU might not work in the US for example
- With FSK modulation: data rate up to 50kbps
- High data rate -> high loss probability
- 2 types of roaming: active and passive
- Classes (A, B or C) define when devices can receive downlink
- Classes can be switched during operation
- Class A most freq used since least power consumed
- Class B can receive periodically using beacons
- Class C can receive continuously -> not suitable for battery operation
- Gateways are stateless
- LoRaWAN networks use Aloha method for communication
- Ethernet or ibackhaul between gateways and network servers
- Network servers take care of: msg consolidation, routing, network control, network and gateway supervision
- ADR: adaptive data rate
- When adding extra gateway, devices will automatically update SF (spreading factor)
- ADR works based on strength: devices closer to gateways use high data rate
- ADR only works for fixed devices, not mobile devices
- The Network Reference Model (NRM) is used to identify the interfaces between elements used in a LoRaWAN network
- RP002 is the companion document to TS001
- 3 factors to consider when deciding on a network:
    - price 
    - gateway access control 
    - privacy
- public networks:
    - verify network coverage before subscription
    - only makes sense if not too many devices
    - cant optimize by adding own gateways
- private networks:
    - for large number of devices 
    - huge overhead: gateways have to be installed and maintained

- The network server is where the LoRaWAN networking protocol terminates, and where data handling takes place
- One can use Semtech's network server for free when prototyping: up to 3 gateways and 10 end devices 
- The join server is used when an end device is joined to a network via Over-the-Air Activation (OTAA) or when a device requests 
  new security session keys via the Rejoin MAC command

- Each end dvice consists of:
    - MCU that runs LoRaWAN stack + power supply
    - peripherals, e.g sensors
    - LoRa Radio
    - Antenna
- Device software layers:
    - drivers
    - middleware
    - application 

- 1/4 wavelength for 868MHz = 8.2cm
- distinction between indoor and outdoor gateways

- LoRaMAC-Node from Semtech: free open-source project
- Also ST's LoRaWAN stack for all STM32Lxx MCUs
- Mbed stack for C++

- LoRaWAN always adds 13 bytes to application payload

- One could, for example, transmit a min|avg|max every five minutes, or you could transmit only when a sensor value changes more than a certain threshold amount. 
  Similarly, you could have transmissions be triggered by motion or other events.
- The data rate should be as fast as possible, to minimize your airtime. 
  Using SF7 and BW125 (SF7BW125) is usually a good place to start.

- FUOTA: firmware update over the air

- Over-the-air-activation (OTAA) vs activation-by-personalization (ABP)
- ABP: device has all the info before connecting to a network + DevEUI + DevAddr
- OTAA: device must issue join request followed by joinaccept from network
- While provisioning devices with ABP, the best practice is to also provision 
  them with all the elements necessary for OTAA (i.e., DevEUI, JoinEUI, and Root Keys)
- OTAA is preferred approach
- All devices must keep a set of persistent values, even after power loss or reset

- Transmission is very power consuming (10 times more than receiving)

- To secure the end devices deployed in the field, crypto chips can be used. Crypto chips, commonly known as secure elements 
  can securely store keys and efficiently encrypt and decrypt data.
- In crypto chips: registers in secure element are write only
- Crypto chips can also be used to verify the authenticity of the software that is uploaded to the device. 
  This mitigates the threat of an attacker overwriting the device firmware
- crypto chip stores network key and derived session keys 
- fast encryption without burdening CPU 
- When using LoRaWAN, ensure that the join server runs in a secure domain that you host or that is hosted by a trusted third party, because this server stores the root keys for the cloud 
  side of the LoRaWAN network. To prevent vendor lock-ins, check to see whether the join server is 
  network-provider agnostic and whether it allows end devices to join a different LoRaWAN network operator over time.
- Out-of-band key provisioning can be a good way to make sure that the people who insert the keys don’t have access to the keys of the devices themselves. 
  It is possible, however, to create an application that uses BLE or NFC to inject the keys directly into the device.
- To mitigate against replay attacks, enable 32-bit frame counters for uplink and downlink messages. Also, use a fixed payload length to mask event-driven 
  activities and help prevent hackers from deriving event-specific information from the payload

- Signal strength testing vs radio planning
- Lookup how to calculate LoRa link budget
- 6dB loss for doubling distance
- FSPL: -150dB/800km
- lookup multipath fading, Fresnel Zone
- 5 to 10 dB of safety margin for link budget
- According to two ray model: -12dB @ 2x distance
- max antenna gain: +2.15dBi (=dipole gain)

- if the internal walls are not brick or concrete, it should be possible to cover most of the 
  building’s interior with a single gateway on the same floor as the end devices.
- You can use radio planning tools to improve the coverage of your network
- There are some tools on the market that have LoRa modules available, such as Forsk Atoll, Infovista Planet, or iBwave
- outdoor gateways are more expensive than indoor due to waterproofness, cavity filters etc
- for indoor: put gateways near glass

- Development stacks:
  - The things stack 
  - ChirpStack (more offgrid promising)
  - zephyr with integrated LoRaWAN stack 
  
- Hardware to buy:
  - nordic power profiler kit II (PPK2)
  - RTL-SDR (software defined radio)
  - For gateway: Dragino LPS8 / MikroTik LR8 / DIY Rasp pi
  - MCU: STM32WL55 (industrial favorite) / Heltech Wifi Lora 32 (ESP-based with OLED) / RakWireless WisBlock (for low-power)

- Software tools:
  - draw.io / Lucidchart
  - Node-Red for data visualization
  - MQTT for debugging
  - octopart for price comparison
  - semtech lora calculator
  - RF line of sight calculator
